\documentclass[11pt,a4paper]{article}

% =========================================================
% FUENTE TIPO ARIAL / HELVETICA
% =========================================================
\usepackage[utf8]{inputenc}
\usepackage[T1]{fontenc}
\usepackage[spanish]{babel}
\usepackage[scaled]{helvet}
\renewcommand{\familydefault}{\sfdefault}

% =========================================================
% PAQUETES
% =========================================================
\usepackage{graphicx}
\usepackage{geometry}
\usepackage{setspace}
\usepackage{fancyhdr}
\usepackage{eso-pic}
\usepackage{titlesec}
\usepackage{array}
\usepackage{tabularx}
\usepackage{multirow}
\usepackage[table]{xcolor}
\usepackage{lastpage}
\usepackage{enumitem}
\usepackage{needspace}
\usepackage{caption}
\usepackage{placeins}
\usepackage[hidelinks]{hyperref}

% =========================================================
% CONFIGURACIÓN GENERAL
% =========================================================
\geometry{
    left=2.35cm,
    right=2.35cm,
    top=3.55cm,
    bottom=4.0cm
}

\setlength{\parindent}{0pt}
\setlength{\parskip}{0.30em}

\setstretch{1.09}
\raggedbottom

% =========================================================
% ESPACIADO DE FIGURAS
% =========================================================
\setlength{\textfloatsep}{8pt}
\setlength{\floatsep}{8pt}
\setlength{\intextsep}{7pt}

% =========================================================
% COLORES CORPORATIVOS
% =========================================================
\definecolor{SIGgris}{RGB}{76,76,76}
\definecolor{SIGgrisClaro}{RGB}{235,235,235}
\definecolor{INGEorange}{RGB}{235,108,54}
\definecolor{AlertaRojo}{RGB}{180,40,40}
\definecolor{AlertaClaro}{RGB}{250,235,235}

% =========================================================
% DATOS DEL DOCUMENTO
% =========================================================
\newcommand{\TituloDocumento}
{GUÍA DE ACCESO REMOTO A PANTALLAS HMI\\
MEDIANTE REALVNC VIEWER}

\newcommand{\Cliente}
{TOMY}

\newcommand{\EmpresaEmisora}
{INGE-LEAN S.A.S.}

% =========================================================
% DATOS DEL ENCABEZADO SIG
% =========================================================
\newcommand{\CodigoDocumento}{SIG-OP-INF-01}
\newcommand{\VersionDocumento}{001}
\newcommand{\FechaDocumento}{30/09/2026}

\newcommand{\TipoDocumento}
{GUÍA DE USUARIO}

\newcommand{\ProcesoDocumento}
{OPERACIONES}

% =========================================================
% FONDO CORPORATIVO
% =========================================================
\IfFileExists{Fondo.jpg}{
    \AddToShipoutPictureBG{
        \includegraphics[
            width=\paperwidth,
            height=\paperheight
        ]{Fondo.jpg}
    }
}{}

% =========================================================
% ENCABEZADO CORPORATIVO SIG
% =========================================================
\newcommand{\EncabezadoSIG}{%
\begingroup

\renewcommand{\arraystretch}{0.90}
\setlength{\tabcolsep}{3pt}

\begin{minipage}{\textwidth}

\begin{tabularx}{\textwidth}
{
|>{\centering\arraybackslash}m{0.255\textwidth}
|>{\centering\arraybackslash}X
|>{\centering\arraybackslash}m{0.085\textwidth}
|>{\centering\arraybackslash}m{0.105\textwidth}|
}

\hline

\multicolumn{4}{|c|}{
    \cellcolor{SIGgris}
    \color{white}
    \rule{0pt}{3.6mm}
    \textbf{\scriptsize SISTEMA INTEGRADO DE GESTIÓN (SIG)}
}
\\

\hline

\multirow{4}{*}{
    \begin{minipage}[c][1.65cm][c]{0.235\textwidth}
        \centering

        \IfFileExists{logo-ingelean.png}{
            \includegraphics[
                width=0.95\linewidth,
                height=1.40cm,
                keepaspectratio
            ]{logo-ingelean.png}
        }{
            {\large\textbf{INGE-LEAN}}
        }

    \end{minipage}
}

&
\multirow{2}{*}{
    \begin{minipage}[c][0.78cm][c]{0.40\textwidth}
        \centering

        \textbf{\scriptsize
        TIPO DE DOCUMENTO: \TipoDocumento
        }

    \end{minipage}
}

&
\cellcolor{SIGgris}
\color{white}
\textbf{\scriptsize Código:}

&
\scriptsize
\CodigoDocumento

\\

\cline{3-4}

&
&
\cellcolor{SIGgris}
\color{white}
\textbf{\scriptsize Versión:}

&
\scriptsize
\VersionDocumento

\\

\cline{2-4}

&
\multirow{2}{*}{
    \begin{minipage}[c][0.78cm][c]{0.40\textwidth}
        \centering

        \textbf{\scriptsize
        PROCESO: \ProcesoDocumento
        }

    \end{minipage}
}

&
\cellcolor{SIGgris}
\color{white}
\textbf{\scriptsize Fecha:}

&
\scriptsize
\FechaDocumento

\\

\cline{3-4}

&
&
\cellcolor{SIGgris}
\color{white}
\textbf{\scriptsize Página:}

&
\scriptsize
\thepage\ de \pageref{LastPage}

\\

\hline

\end{tabularx}

\end{minipage}

\endgroup
}

% =========================================================
% CONFIGURACIÓN DEL ENCABEZADO
% =========================================================
\pagestyle{fancy}
\fancyhf{}

\setlength{\headheight}{75pt}
\setlength{\headsep}{10pt}

\renewcommand{\headrulewidth}{0pt}
\renewcommand{\footrulewidth}{0pt}

\fancyhead[C]{\vspace{-1.5cm}\EncabezadoSIG}

% =========================================================
% FORMATO DE SECCIONES
% =========================================================
\titleformat{\section}
{\normalfont\large\bfseries}
{\thesection.}
{0.6em}
{}

\titleformat{\subsection}
{\normalfont\normalsize\bfseries}
{\thesubsection.}
{0.5em}
{}

\titlespacing*{\section}
{0pt}
{0.95em}
{0.38em}

\titlespacing*{\subsection}
{0pt}
{0.75em}
{0.28em}

% =========================================================
% LISTAS
% =========================================================
\setlist[itemize]{
    leftmargin=1.60em,
    itemsep=0.12em,
    topsep=0.16em,
    parsep=0pt
}

\setlist[enumerate]{
    leftmargin=1.85em,
    itemsep=0.12em,
    topsep=0.16em,
    parsep=0pt
}

% =========================================================
% CONFIGURACIÓN DE PIES DE IMAGEN
% =========================================================
\captionsetup{
    font=small,
    labelfont=bf,
    justification=centering,
    singlelinecheck=false,
    skip=4pt
}

% =========================================================
% DOCUMENTO
% =========================================================
\begin{document}

% =========================================================
% PORTADA
% =========================================================

\vspace*{0.02cm}

\begin{center}

{\LARGE\bfseries
\TituloDocumento
}

\vspace{0.55cm}

\vspace{0.55cm}

\begin{tabular}{>{\bfseries}l l}

Cliente: &
\Cliente
\\[0.14cm]

Empresa ejecutora: &
\EmpresaEmisora
\\[0.14cm]

Aplicación: &
Acceso remoto a pantallas HMI
\\

\end{tabular}

\end{center}

\vspace{0.30cm}

% =========================================================
% 1. OBJETIVO
% =========================================================
\section{Objetivo}

El presente documento tiene como finalidad explicar el procedimiento
básico para acceder de forma remota a las pantallas HMI mediante el
software RealVNC Viewer.

La guía aplica a las pantallas HMI correspondientes a los equipos Horno,
Laminación y Hojadrador conectados a la red interna.

Además del procedimiento de conexión, se establecen las consideraciones
de seguridad y responsabilidad asociadas al uso del acceso remoto.

% =========================================================
% 2. ADVERTENCIA IMPORTANTE
% =========================================================
\section{Advertencia sobre el acceso remoto}

\begin{center}

\renewcommand{\arraystretch}{1.18}

\begin{tabularx}{0.95\textwidth}{
|>{\raggedright\arraybackslash}X|
}

\hline

\cellcolor{AlertaClaro}

\textbf{\color{AlertaRojo}IMPORTANTE:}

El acceso mediante RealVNC Viewer no corresponde únicamente a una
visualización de la pantalla.

Al establecer la conexión remota, el usuario puede interactuar
directamente con la HMI de la misma manera que si estuviera operándola
físicamente frente al equipo.

\\

\hline

\end{tabularx}

\end{center}

Cualquier comando realizado desde la conexión remota puede generar una
acción directa sobre el sistema.

Dependiendo de la pantalla y del modo de operación disponible, el usuario
puede tener acceso a funciones como:

\begin{itemize}

    \item accionamiento o detención de equipos desde los mandos manuales;

    \item modificación de frecuencias o referencias de operación;

    \item cambio de parámetros utilizados durante el proceso;

    \item activación o desactivación de elementos disponibles desde la HMI;

    \item ejecución de comandos mientras otro operador se encuentra
    trabajando directamente sobre la máquina.

\end{itemize}

Un cambio incorrecto de parámetros o una maniobra realizada sin verificar
previamente el estado del equipo puede afectar la operación de la máquina
o el proceso productivo.

Por esta razón, el acceso remoto debe ser utilizado únicamente por
personal autorizado y con conocimiento de la operación del equipo.

\textbf{Las acciones realizadas durante una sesión remota y las
consecuencias derivadas de los comandos o modificaciones ejecutadas desde
la HMI serán responsabilidad de la persona que realice la conexión.}

Antes de realizar cualquier maniobra se recomienda confirmar con el
personal de operación el estado actual de la máquina.

% =========================================================
% 3. APERTURA DE REALVNC
% =========================================================

\Needspace{0.54\textheight}

\section{Apertura de RealVNC Viewer}

Al abrir RealVNC Viewer se mostrará la pantalla principal del programa.

Si todavía no se han guardado conexiones, el área central aparecerá
vacía y no se mostrarán dispositivos registrados.

En la parte superior de la ventana se encuentra la barra utilizada para
ingresar la dirección IP del equipo al que se desea acceder.

\vspace{0.10cm}

\begin{center}

\IfFileExists{Pantalla_Incial.jpeg}{

    \includegraphics[
        width=0.93\textwidth,
        height=0.39\textheight,
        keepaspectratio
    ]{Pantalla_Incial.jpeg}

}{
    \fbox{
        \parbox[c][6.2cm][c]{0.88\textwidth}{
            \centering
            Pantalla inicial de RealVNC Viewer
        }
    }
}

\captionof{figure}{
Pantalla inicial de RealVNC Viewer antes de registrar las conexiones.
}

\end{center}

% =========================================================
% 4. INGRESO DE DIRECCIÓN IP
% =========================================================

\Needspace{0.72\textheight}

\section{Ingreso de la dirección IP}

Para iniciar una nueva conexión, se debe ingresar en la barra superior la
dirección IP correspondiente a la pantalla HMI que se desea controlar.

En la siguiente imagen se resalta el campo donde debe escribirse la
dirección IP.

\vspace{0.10cm}

\begin{center}

\IfFileExists{Ubicacion_IP.jpeg}{

    \includegraphics[
        width=0.96\textwidth,
        height=0.39\textheight,
        keepaspectratio
    ]{Ubicacion_IP.jpeg}

}{
    \fbox{
        \parbox[c][6.0cm][c]{0.90\textwidth}{
            \centering
            Ubicación para ingresar la dirección IP
        }
    }
}

\captionof{figure}{
Campo de RealVNC Viewer utilizado para ingresar la dirección IP de la HMI.
}

\end{center}

\vspace{0.05cm}

Las direcciones configuradas son las siguientes:

\vspace{0.08cm}

\begin{center}

\renewcommand{\arraystretch}{1.18}

\begin{tabularx}{0.76\textwidth}{
|>{\raggedright\arraybackslash}X
|>{\centering\arraybackslash}m{4.0cm}|
}

\hline

\cellcolor{SIGgrisClaro}
\textbf{Pantalla HMI}
&
\cellcolor{SIGgrisClaro}
\textbf{Dirección IP}
\\

\hline

HMI\_HORNO
&
192.168.4.81
\\

\hline

HMI\_LAMINACION
&
192.168.4.90
\\

\hline

HMI\_HOJADRADOR
&
192.168.4.109
\\

\hline

\end{tabularx}

\end{center}

\vspace{0.02cm}

Una vez ingresada la dirección IP correspondiente, se debe presionar
\textbf{Enter} para iniciar la conexión.

% =========================================================
% 5. CONTRASEÑA
% =========================================================

\Needspace{9\baselineskip}

\section{Contraseña de acceso}

Durante la primera conexión, RealVNC Viewer solicitará la contraseña
configurada para el acceso remoto a la HMI.

La contraseña actualmente establecida es:

\vspace{-0.08cm}

\begin{center}

\fbox{
    \parbox[c][0.75cm][c]{4.6cm}{
        \centering
        \large\textbf{12345678}
    }
}

\end{center}

\vspace{-0.08cm}

Después de ingresar correctamente la contraseña, se podrá establecer la
conexión con la pantalla seleccionada.

La contraseña debe ser utilizada únicamente por personal autorizado y no
debe compartirse con personas ajenas a la operación o mantenimiento de
los equipos.

% =========================================================
% 6. CONEXIONES GUARDADAS
% =========================================================

\Needspace{0.66\textheight}

\section{Conexiones guardadas}

Una vez realizada la conexión, RealVNC Viewer permite conservar los
dispositivos registrados para facilitar accesos posteriores.

Después de guardar las conexiones, las tres pantallas aparecerán en la
ventana principal del programa.

\vspace{0.10cm}

\begin{center}

\IfFileExists{Pantallas_Guardadas.png}{

    \includegraphics[
        width=0.94\textwidth,
        height=0.39\textheight,
        keepaspectratio
    ]{Pantallas_Guardadas.png}

}{
    \fbox{
        \parbox[c][6.0cm][c]{0.88\textwidth}{
            \centering
            Pantallas HMI guardadas en RealVNC Viewer
        }
    }
}

\captionof{figure}{
Pantallas HMI registradas y disponibles para acceso desde RealVNC Viewer.
}

\end{center}

\vspace{0.02cm}

Las conexiones corresponden a:

\begin{itemize}

    \item \textbf{192.168.4.81} -- HMI\_HORNO;

    \item \textbf{192.168.4.90} -- HMI\_LAMINACION;

    \item \textbf{192.168.4.109} -- HMI\_HOJADRADOR.

\end{itemize}

Mientras el computador se encuentre conectado al mismo segmento de red
de las pantallas HMI, para ingresar nuevamente será suficiente seleccionar
la pantalla correspondiente y abrir la conexión.

% =========================================================
% 7. CONDICIONES DE CONEXIÓN
% =========================================================

\Needspace{14\baselineskip}

\section{Condiciones necesarias para el acceso}

Para establecer correctamente una conexión remota se debe verificar que:

\begin{itemize}

    \item el computador se encuentre conectado a la red correspondiente;

    \item el computador y la HMI se encuentren dentro del mismo segmento
    de red;

    \item exista comunicación con la dirección IP de la HMI;

    \item la pantalla HMI se encuentre energizada y disponible;

    \item el usuario cuente con autorización para realizar el acceso.

\end{itemize}

Si alguna de estas condiciones no se cumple, la conexión puede no estar
disponible aunque el dispositivo haya sido registrado previamente en
RealVNC Viewer.

% =========================================================
% 8. CONSIDERACIONES DE USO
% =========================================================

\Needspace{14\baselineskip}

\section{Consideraciones de uso}

Durante una sesión remota, las acciones realizadas se reflejan
directamente sobre la HMI del equipo.

Antes de ejecutar comandos manuales o modificar valores de operación se
debe verificar el estado de la máquina y coordinar la intervención con el
personal responsable del proceso.

No se deben modificar frecuencias, referencias, parámetros de control ni
mandos manuales si no se conoce previamente el efecto que dicho cambio
puede generar sobre el sistema.

Cuando exista personal trabajando directamente sobre la máquina, deberá
existir coordinación previa antes de realizar cualquier operación desde
RealVNC.

Al finalizar una revisión o intervención remota, se recomienda cerrar la
sesión para evitar operaciones accidentales posteriores.

% =========================================================
% 9. CONCLUSIÓN
% =========================================================

\Needspace{12\baselineskip}

\section{Conclusión}

RealVNC Viewer permite acceder remotamente a las pantallas HMI de Horno,
Laminación y Hojadrador desde un computador conectado a la misma red.

Esta herramienta facilita actividades de revisión y soporte sin necesidad
de encontrarse físicamente frente a cada pantalla.

Sin embargo, el acceso remoto permite interactuar directamente con los
controles disponibles en la HMI y debe utilizarse con el mismo nivel de
precaución que una operación realizada localmente.

\textbf{El usuario que establezca la conexión será responsable de las
acciones ejecutadas durante la sesión remota y deberá evitar cualquier
modificación que pueda afectar la operación normal del equipo o del
proceso productivo.}

\end{document}